%% ============================================================
\section{System Architecture}
\label{sec:architecture}
%% ============================================================

\subsection{Architectural Philosophy}

\echoray{} is designed around three guiding principles:

\begin{enumerate}[leftmargin=*, label=\textbf{P\arabic*.}]
  \item \textbf{Separation of Concerns.} Business logic, data access, and
    transport are cleanly layered: Routes $\rightarrow$ Controllers
    $\rightarrow$ Services $\rightarrow$ Prisma ORM, with Socket.io
    orthogonal to the REST layer.
  \item \textbf{Type Safety End-to-End.} Both client and server are written in
    TypeScript; Prisma generates typed query methods from the schema; Zod
    validates and narrows all external inputs at runtime.
  \item \textbf{Cloud-Native Statelessness.} Application servers are
    horizontally scalable; the only stateful services are the PostgreSQL
    database and Redis cache, both of which are externally managed.
\end{enumerate}

\subsection{Layered Component Model}

\Cref{fig:arch_layers} depicts the full layered architecture.

\begin{figure}[H]
  \centering
  \begin{tikzpicture}[
    font=\small,
    layer/.style={
      rectangle, rounded corners=3pt,
      minimum width=11cm, minimum height=0.9cm,
      text centered, draw, thick
    },
    lbl/.style={font=\bfseries\footnotesize, anchor=west},
  ]
    \node[layer, fill=blue!10]  (l1) at (0,0)    {Browser: Next.js 15 + React 19 + Monaco Editor + Redux + TailwindCSS};
    \node[layer, fill=cyan!10]  (l2) at (0,-1.1) {Transport: Socket.io 4.8 (WebSocket) \quad|\quad Axios HTTP/REST};
    \node[layer, fill=green!10] (l3) at (0,-2.2) {API Gateway: Express 5 — Routes / Controllers / Zod Validation / Helmet};
    \node[layer, fill=yellow!15](l4) at (0,-3.3) {Business Logic: Services (User, Project, AI, Redis)};
    \node[layer, fill=orange!15](l5) at (0,-4.4) {Data Access: Prisma ORM with Accelerate \quad|\quad Redis (ioredis)};
    \node[layer, fill=red!10]   (l6) at (0,-5.5) {Persistence: PostgreSQL (Supabase/Neon) \quad|\quad Redis Cloud / Upstash};
    \node[layer, fill=purple!10](l7) at (0,-6.6) {External AI: Google Gemini 2.5 Flash (REST via \texttt{@google/generative-ai})};

    \foreach \a/\b in {l1/l2,l2/l3,l3/l4,l4/l5,l5/l6,l6/l7}{
      \draw[-Stealth,thick] (\a.south) -- (\b.north);
      \draw[-Stealth,thick] (\b.north) -- (\a.south);
    }
  \end{tikzpicture}
  \caption{Layered architecture of \echoray. Bidirectional arrows denote
    request/response or event flow between adjacent layers.}
  \label{fig:arch_layers}
\end{figure}

\subsection{Database Schema}

The data model is deliberately minimal, exposing two principal entities:
\texttt{User} and \texttt{Project}, linked by a many-to-many association
representing project membership and a one-to-many association representing
project leadership.

\begin{lstlisting}[language=SQL,caption={Prisma schema (\texttt{schema.prisma}). The
  \texttt{fileTree} column stores the entire project file structure as a
  JSON blob, enabling schema-free project content evolution.},
  label=lst:schema]
model User {
  id              Int       @id @default(autoincrement())
  email           String    @unique
  password        String
  createdAt       DateTime  @default(now())
  projects        Project[] @relation("UserProjects")
  leadingProjects Project[] @relation("ProjectLeader")
}

model Project {
  id        Int      @id @default(autoincrement())
  name      String   @unique
  users     User[]   @relation("UserProjects")
  fileTree  Json     @default("{}")
  createdAt DateTime @default(now())
  leaderId  Int
  leader    User     @relation("ProjectLeader",
                      fields: [leaderId], references: [id])
}
\end{lstlisting}

The \texttt{fileTree} JSON column stores the project's complete virtual file
system as a flat key-value map of the form:

\begin{equation}
  \mathcal{F} : \text{path} \to \bigl\{\ \texttt{file}:
  \{\,\texttt{contents}: \text{string}\,\}\ \bigr\}
  \label{eq:filetree}
\end{equation}

This representation is directly compatible with the WebContainer
API~\cite{stackblitz2023wc} file-system mounting interface, enabling
in-browser execution without any transformation.

\subsection{REST API Surface}

\echoray{} exposes sixteen REST endpoints across three route namespaces,
as summarised in \cref{tab:api}.

\begin{table}[H]
  \centering
  \caption{REST API endpoint catalogue (\texttt{server/src/routes/}).
    Protected endpoints require a valid JWT in the \texttt{Authorization}
    header or \texttt{token} cookie; the token is also verified against the
    Redis blacklist.}
  \label{tab:api}
  \setlength{\tabcolsep}{4pt}
  \begin{tabular}{llll}
    \toprule
    \textbf{Method} & \textbf{Endpoint} & \textbf{Auth} & \textbf{Purpose} \\
    \midrule
    POST   & \texttt{/api/v1/user/register}        & No  & Create user account \\
    POST   & \texttt{/api/v1/user/login}            & No  & Authenticate, issue JWT \\
    GET    & \texttt{/api/v1/user/profile}          & Yes & Fetch own profile \\
    GET    & \texttt{/api/v1/user/logout}           & Yes & Revoke token (Redis) \\
    GET    & \texttt{/api/v1/user/getAll}           & Yes & List all other users \\
    GET    & \texttt{/api/v1/user/getUserByEmail}   & Yes & Find user by e-mail \\
    \midrule
    POST   & \texttt{/api/v1/project/create}        & Yes & Create project \\
    GET    & \texttt{/api/v1/project/getAll}        & Yes & List user's projects \\
    POST   & \texttt{/api/v1/project/addUsers}      & Yes & Add collaborators (leader) \\
    GET    & \texttt{/api/v1/project/get-project/:id}& Yes & Get project details \\
    DELETE & \texttt{/api/v1/project/delete/:id}   & Yes & Delete project (leader) \\
    POST   & \texttt{/api/v1/project/removeUsers}  & Yes & Remove collaborators (leader) \\
    PUT    & \texttt{/api/v1/project/update-file-tree}& Yes & Persist file tree \\
    \midrule
    GET    & \texttt{/api/v1/ai/get-result}         & No  & Query Gemini (prompt) \\
    \bottomrule
  \end{tabular}
\end{table}

\subsection{Deployment Topology}

\Cref{fig:deploy} depicts the production deployment topology. The front end
is served as a hybrid SSR/SSG Next.js application from Vercel's edge network.
The back end runs on Render or Railway as a persistent Node.js process
listening on port~5000. The database is hosted on Supabase or Neon (both
provide serverless PostgreSQL with connection pooling), accessed via Prisma
Accelerate for query-level caching. Redis is provisioned on Upstash or Redis
Cloud, providing a geo-distributed key-value store for token blacklisting.

\begin{figure}[H]
  \centering
  \begin{tikzpicture}[
    font=\small,
    svc/.style={
      rectangle, rounded corners=4pt,
      minimum width=2.5cm, minimum height=0.8cm,
      text centered, align=center, draw, thick
    },
    arr/.style={-Stealth, thick},
  ]
    \node[svc, fill=blue!15]   (browser) at (0,0)     {Browser\\(Vercel CDN)};
    \node[svc, fill=green!15]  (api)     at (4,0)     {Express API\\(Render/Railway)};
    \node[svc, fill=orange!15] (pg)      at (8,1)     {PostgreSQL\\(Supabase/Neon)};
    \node[svc, fill=red!15]    (redis)   at (8,-1)    {Redis\\(Upstash)};
    \node[svc, fill=purple!15] (gemini)  at (4,-2)    {Google Gemini\\2.5 Flash};

    \draw[arr, bend left=10]  (browser.east)  to node[above]{\tiny HTTPS/WSS} (api.west);
    \draw[arr, bend left=10]  (api.west)      to (browser.east);
    \draw[arr]                (api.east)       to[bend left=15]  (pg.west);
    \draw[arr]                (pg.west)        to[bend left=15]  (api.east);
    \draw[arr]                (api.east)       to[bend right=15] (redis.west);
    \draw[arr]                (redis.west)     to[bend right=15] (api.east);
    \draw[arr]                (api.south)      to (gemini.north);
    \draw[arr]                (gemini.north)   to (api.south);
  \end{tikzpicture}
  \caption{Production deployment topology of \echoray. Browser clients
    communicate with the Express API over HTTPS for REST calls and WSS for
    Socket.io. The API server interacts with PostgreSQL via Prisma, Redis via
    ioredis, and Google Gemini via the official Node.js SDK.}
  \label{fig:deploy}
\end{figure}

\subsection{Key Design Decisions}

\textbf{JSON-blob file storage.}
Storing the entire file tree as a single \texttt{Json} column trades
fine-grained SQL access for implementation simplicity and atomic updates.
Because projects in \echoray{} are small-to-medium prototypes (comparable in
scale to StackBlitz sandboxes), this approach is practical; migration to a
normalised file-entity model is deferred to future work (\cref{sec:conclusion}).

\textbf{Express~5 over Next.js API routes.}
The API is hosted in a separate Express server rather than Next.js API routes
to allow independent scaling, easier Socket.io integration (which requires a
persistent Node.js process), and deployment to general-purpose PaaS providers.

\textbf{Prisma Accelerate.}
Prisma Accelerate provides a connection pool and an LRU query cache at the
database edge, reducing median query latency from~$\sim$30\,ms (direct
connection) to~$\sim$8\,ms for read-heavy endpoints (see \cref{sec:evaluation}).

\textbf{TypeScript end-to-end.}
TypeScript catches a significant class of type errors at compile time. Zod
schemas in \texttt{utils/schema.ts} add a runtime validation layer, ensuring
that even valid TypeScript types that carry semantically invalid values (e.g.,
an empty string as a project name) are rejected before reaching the database.
